\setcounter{section}{3}

  \zwischenueberschrift{Soal latihan}

\inputaufgabe
{}
{

Misalkan
\maabb {\gamma} { I } {\R^2
} {}
sebuah
\definitionsverweis {kurva berparametrisasi panjang busur}{}{}
yang dua kali terdiferensialkan secara kontinu, dan misalkan
\maabbeledisp {\delta} { -I } {\R^2
} { s } { \delta(s) = \gamma( -s)
} {,}
yaitu kurva yang ditelusuri dalam arah sebaliknya. Buktikan bahwa
\definitionsverweis {kelengkungan}{}{}
$\delta$ di $s$ merupakan negatif dari kelengkungan
\mathl{\gamma}{} di $-s$.

}
{} {}

\inputaufgabe
{}
{

Misalkan
\maabb {\gamma} { I } {\R^2
} {}
sebuah
\definitionsverweis {kurva berparametrisasi panjang busur}{}{}
yang dua kali terdiferensialkan secara kontinu, dan
\maabbdisp {L} {\R^2} {\R^2
} {}
sebuah
\definitionsverweis {rotasi}{}{.}
Buktikan bahwa
\definitionsverweis {kelengkungan}{}{}
$\gamma$ sama dengan kelengkungan
\mathl{L\circ \gamma}{}.

}
{} {}

\inputaufgabe
{}
{

Buktikan
Lema 3.6
dalam kasus dengan orientasi nonstandar.

}
{} {}

\inputaufgabe
{}
{

Misalkan
\maabb {\gamma} { I } {\R^2
} {}
sebuah
\definitionsverweis {kurva berparametrisasi panjang busur}{}{}
yang dua kali terdiferensialkan secara kontinu, dan
\maabbdisp {L} {\R^2} {\R^2
} {}
sebuah
\definitionsverweis {homoteti skalar}{}{}
dengan faktor homoteti

\mavergleichskette
{\vergleichskette
{ s
}
{ \neq }{ 0
}
{  }{
}
{    }{
}
{   }{
}
}
{}{}{.}
Bagaimana hubungan antara
\definitionsverweis {kelengkungan}{}{}
$\gamma$ dan kelengkungan
\mathl{L\circ \gamma}{?}

}
{} {}

\inputaufgabe
{}
{

Misalkan
\maabbeledisp {\gamma} {\R} { \R^2
} { t } { \begin{pmatrix}  \cos \left(  \omega t \right)  \\ \sin  \left(  \omega t \right)    \end{pmatrix}
} {,}
dengan

\mavergleichskette
{\vergleichskette
{ \omega
}
{ > }{ 0
}
{  }{
}
{    }{
}
{   }{
}
}
{}{}{.}
Buktikan bahwa terdapat tak berhingga banyak gerak melingkar dengan norma kecepatan konstan yang memiliki nilai dan turunan pertama yang sama dengan $\gamma$ pada

\mavergleichskette
{\vergleichskette
{ t
}
{ = }{ 0
}
{  }{
}
{    }{
}
{   }{
}
}
{}{}{.}

}
{} {}

Soal berikut penting dalam pembuatan komidi putar. Kesenangan menaiki komidi putar berasal dari percepatan, bukan dari kecepatan.

\inputaufgabe
{}
{

Misalkan
\maabbeledisp {\gamma} {\R} { \R^2
} { t } { \begin{pmatrix}  \cos \left(  \omega t \right)  \\ \sin  \left(  \omega t \right)    \end{pmatrix}
} {,}
dengan

\mavergleichskette
{\vergleichskette
{ \omega
}
{ > }{ 0
}
{  }{
}
{    }{
}
{   }{
}
}
{}{}{.}
Buktikan bahwa terdapat tak berhingga banyak gerak melingkar dengan norma kecepatan konstan yang berimpit dengan $\gamma$ pada

\mavergleichskette
{\vergleichskette
{ t
}
{ = }{ 0
}
{  }{
}
{    }{
}
{   }{
}
}
{}{}{}
dan juga memiliki percepatan yang sama di sana.

}
{} {}

\inputaufgabegibtloesung
{}
{

Misalkan
\maabbeledisp {\gamma} {\R} { \R^2
} { t } { \begin{pmatrix}  \cos \left(  \omega t \right)  \\ \sin  \left(  \omega t \right)    \end{pmatrix}
} {,}
dengan

\mavergleichskette
{\vergleichskette
{ \omega
}
{ > }{ 0
}
{  }{
}
{    }{
}
{   }{
}
}
{}{}{.}
Buktikan bahwa tidak ada gerak melingkar lain dengan norma kecepatan konstan yang memiliki nilai serta turunan pertama dan kedua yang sama dengan $\gamma$ pada

\mavergleichskette
{\vergleichskette
{ t
}
{ = }{ 0
}
{  }{
}
{    }{
}
{   }{
}
}
{}{}{.}

}
{} {}

\inputaufgabe
{}
{

Misalkan

\mavergleichskette
{\vergleichskette
{ W
}
{ \subseteq }{ \R^2
}
{  }{
}
{    }{
}
{   }{
}
}
{}{}{}
\definitionsverweis {terbuka}{}{,}
\maabb {h} { W } { \R
} {}
sebuah fungsi yang dua kali
\definitionsverweis {terdiferensialkan secara kontinu}{}{,}
dan

\mavergleichskette
{\vergleichskette
{ Y
}
{ = }{ h^{-1}(c)
}
{  }{
}
{    }{
}
{   }{
}
}
{}{}{}
adalah
\definitionsverweis {serat}{}{}
di atas

\mavergleichskette
{\vergleichskette
{ c
}
{ \in }{ \R
}
{  }{
}
{    }{
}
{   }{
}
}
{}{}{,}
dengan $h$
\definitionsverweis {reguler}{}{}
di setiap titik pada $Y$. Buktikan bahwa
\maabbdisp {\gamma} { I } { Y
} {}
merupakan
\definitionsverweis {kurva geodesik}{}{}
jika dan hanya jika $\gamma$
\definitionsverweis {berparametrisasi panjang busur}{}{.}

}
{} {}

\inputaufgabe
{}
{

\bild{ \begin{center}
\includegraphics[width=5.5cm]{ \bildeinlesungsvg {Euler spiral} {svg} }
\end{center}
\bildtext {} }

\bildlizenz { Euler spiral.svg } {} {AdiJapan} {Commons} {CC-by-sa 3.0} {}

Kita tinjau kurva terdiferensialkan
\zusatzklammer {yang disebut \stichwort {klotoid} {}} {} {}
\maabbeledisp {} {\R} {\R^2
} { t } { \gamma(t)
} {,}
dengan

\mavergleichskettedisp
{\vergleichskette
{ \gamma(t)
}
{ =} { \sqrt{\pi} \int_0^t  \begin{pmatrix}  \cos  { \frac{   \pi s^2 }{  2 } }  \\ \sin  { \frac{   \pi s^2 }{  2 } }     \end{pmatrix}  ds
}
{ } {
}
{ } {
}
{ } {
}
}
{}{}{.}
Buktikan bahwa
\definitionsverweis {kelengkungan}{}{}
kurva ini memenuhi persamaan

\mavergleichskettedisp
{\vergleichskette
{ \kappa(t)
}
{ =} { \sqrt{\pi} t
}
{ } {
}
{ } {
}
{ } {
}
}
{}{}{.}

}
{} {}

\inputaufgabe
{}
{

Misalkan
\maabb {f} { I } {\R
} {}
sebuah fungsi yang dua kali
\definitionsverweis {terdiferensialkan secara kontinu}{}{.}
Tentukan
\definitionsverweis {kelengkungan}{}{}
dan lingkaran kelengkungan dari grafik yang terkait.

}
{} {}

\inputaufgabe
{}
{

Misalkan
\maabb {f} {\R} {\R
} {}
dua kali
\definitionsverweis {terdiferensialkan secara kontinu}{}{}
dan

\mavergleichskettedisp
{\vergleichskette
{ \alpha(t)
}
{ \defeq} { \begin{pmatrix}  \cos f(t)  \\ \sin f(t)   \end{pmatrix}
}
{ } {
}
{ } {
}
{ } {
}
}
{}{}{.}
Tentukan
\definitionsverweis {kelengkungan}{}{}
$\alpha$ dan
\definitionsverweis {lingkaran kelengkungan}{}{}
untuk

\mavergleichskette
{\vergleichskette
{ t
}
{ \in }{ \R
}
{  }{
}
{    }{
}
{   }{
}
}
{}{}{}
dengan syarat $f'(t) \neq 0$, menggunakan rumus-rumus dalam
Lema 3.10.

}
{} {}

\inputaufgabe
{}
{

Tentukan
\definitionsverweis {kelengkungan}{}{}
\definitionsverweis {grafik}{}{}
\maabbdisp {\sin} {\R} {\R
} {}
untuk

\mavergleichskette
{\vergleichskette
{ t
}
{ = }{ { \frac{   \pi }{  2 } }
}
{  }{
}
{    }{
}
{   }{
}
}
{}{}{.}

}
{} {}

\inputaufgabe
{}
{

Tentukan
\definitionsverweis {kelengkungan}{}{}
spiral logaritmik
\maabbeledisp {\gamma} {\R} { \R^2
} { t } { \gamma(t) = e^t \left(  \cos  t   , \,   \sin  t                   \right)
} {,}
untuk setiap

\mavergleichskette
{\vergleichskette
{ t
}
{ \in }{ \R
}
{  }{
}
{    }{
}
{   }{
}
}
{}{}{.}

}
{} {}

Sebuah fungsi
\maabb {f} { I } {\R
} {}
pada interval terbuka

\mavergleichskette
{\vergleichskette
{ I
}
{ \subseteq }{ \R
}
{  }{
}
{    }{
}
{   }{
}
}
{}{}{}
disebut
\definitionswort {analitik}{,}
jika pada setiap titik

\mavergleichskette
{\vergleichskette
{ a
}
{ \in }{ I
}
{  }{
}
{    }{
}
{   }{
}
}
{}{}{}
fungsi tersebut dapat dinyatakan oleh suatu
\definitionsverweis {deret pangkat}{}{.}

Sebuah kurva disebut analitik jika setiap fungsi komponennya analitik.

\inputaufgabe
{}
{

Misalkan
\maabb {\kappa} { I } {\R
} {}
sebuah
\definitionsverweis {fungsi analitik}{}{.}
Buktikan bahwa terdapat kurva analitik
\maabb {\gamma} { I } {\R^2
} {}
yang
\definitionsverweis {kelengkungan}{}{}
di titik $\gamma(t)$ sama dengan $\kappa(t)$.

}
{} {}

\inputaufgabe
{}
{

\bild{ \begin{center}
\includegraphics[width=5.5cm]{ \bildeinlesungsvg {Evolute-parab} {svg} }
\end{center}
\bildtext {} }

\bildlizenz { Evolute-parab.svg } {} {Ag2gaeh} {Commons} {CC-by-sa 4.0} {}

Tentukan
\definitionsverweis {evolut}{}{}
kurva
\maabbeledisp {} {\R} {\R^2
} { t } { \begin{pmatrix}  t  \\ t^2    \end{pmatrix}
} {.}

}
{} {}

\inputaufgabegibtloesung
{}
{

Tentukan
\definitionsverweis {kelengkungan}{}{}
di setiap titik lingkaran satuan yang diberikan secara implisit oleh

\mavergleichskettedisp
{\vergleichskette
{ h(x,y)
}
{ =} { x^2+y^2
}
{ =} { 1
}
{ } {
}
{ } {
}
}
{}{}{}
menggunakan
Lema 3.11
dan medan gradien $h$.

}
{} {}

  \zwischenueberschrift{Soal untuk dikumpulkan}

\inputaufgabe
{2}
{

Tentukan
\definitionsverweis {kelengkungan}{}{}
\definitionsverweis {grafik}{}{}
fungsi eksponensial
\maabbdisp {\exp} {\R} {\R
} {}
untuk setiap

\mavergleichskette
{\vergleichskette
{ t
}
{ \in }{ \R
}
{  }{
}
{    }{
}
{   }{
}
}
{}{}{.}

}
{} {}

\inputaufgabe
{2}
{

Tentukan
\definitionsverweis {kelengkungan}{}{}
spiral Archimedes
\maabbeledisp {\gamma} {\R} { \R^2
} { t } { \gamma(t) = t \left(  \cos  t   , \,   \sin  t                   \right)
} {,}
untuk setiap

\mavergleichskette
{\vergleichskette
{ t
}
{ \in }{ \R
}
{  }{
}
{    }{
}
{   }{
}
}
{}{}{.}

}
{} {}

\inputaufgabe
{4}
{

Berikan sebuah contoh kurva berparametrisasi panjang busur
\maabbdisp {\gamma} {\R} {\R^2
} {}
yang dua kali terdiferensialkan secara kontinu, sedemikian sehingga untuk dua waktu

\mavergleichskette
{\vergleichskette
{ t_0
}
{ \neq }{ t_1
}
{  }{
}
{    }{
}
{   }{
}
}
{}{}{}
berlaku kesamaan

\mavergleichskette
{\vergleichskette
{ \gamma (t_0)
}
{ = }{ \gamma(t_1)
}
{  }{
}
{    }{
}
{   }{
}
}
{}{}{}
dan
\definitionsverweis {lingkaran-lingkaran kelengkungan}{}{}
pada
\mathkor {} {t_0} {dan} {t_1} {}
tidak sama.

}
{} {}

\inputaufgabe
{4}
{

Tentukan
\definitionsverweis {evolut}{}{}
kurva
\maabbeledisp {} {\R} {\R^2
} { t } { \begin{pmatrix}  2 \cos  t  \\ 3 \sin  t    \end{pmatrix}
} {.}
Di titik-titik manakah turunan evolut mempunyai titik nol?

}
{} {}

\inputaufgabe
{4}
{

Tentukan
\definitionsverweis {kelengkungan}{}{}
lingkaran satuan yang diberikan secara implisit oleh

\mavergleichskettedisp
{\vergleichskette
{ h(x,y)
}
{ =} { x^4+y^4
}
{ =} { 1
}
{ } {
}
{ } {
}
}
{}{}{}
menggunakan
Lema 3.11
dan medan gradien $h$ di titik-titik berikut.
\aufzaehlungzwei {
\mavergleichskette
{\vergleichskette
{ P
}
{ = }{ \begin{pmatrix}  1  \\ 0   \end{pmatrix}
}
{  }{
}
{    }{
}
{   }{
}
}
{}{}{,}
} {
\mavergleichskette
{\vergleichskette
{ Q
}
{ = }{ \begin{pmatrix}  { \frac{   1  }{  \sqrt[4]{2}  } }  \\ { \frac{   1  }{  \sqrt[4]{2}  } }     \end{pmatrix}
}
{  }{
}
{    }{
}
{   }{
}
}
{}{}{.}
}

}
{} {}
